\documentclass[11pt]{article}

\usepackage[utf8]{inputenc}
\usepackage[T1]{fontenc}
\usepackage{lmodern}
\usepackage{geometry}
\usepackage{amsmath,amssymb,amsfonts,amsthm,mathtools}
\usepackage{bm}
\usepackage{enumitem}
\usepackage{microtype}
\usepackage{hyperref}
\usepackage{titlesec}
\usepackage{booktabs}
\usepackage{array}
\usepackage{setspace}
\usepackage{mathrsfs}
\usepackage{physics}

\geometry{margin=1in}
\hypersetup{colorlinks=true, linkcolor=black, urlcolor=blue, citecolor=black}
\setlist[enumerate]{topsep=0.4em,itemsep=0.35em}
\setlist[itemize]{topsep=0.3em,itemsep=0.25em}
\titleformat{\section}{\large\bfseries}{\thesection.}{0.5em}{}
\titleformat{\subsection}{\normalsize\bfseries}{\thesubsection.}{0.5em}{}
\setstretch{1.06}

\newcommand{\Aether}{\AE{}ther}
\newcommand{\Aflow}{\AE{}ther-flow}
\newcommand{\TheoryName}{\Aflow{} Interpretation of Relativity}
\newcommand{\TheoryProgram}{\Aether{} / \Aflow{} framework}
\newcommand{\Sub}{\mathcal{A}}
\newcommand{\BridgeMetric}{\mathfrak g}
\newcommand{\ObserverMap}{\Xi}
\newcommand{\ProjSub}{P}
\newcommand{\SpatialD}{\mathcal D}
\newcommand{\LapPerp}{\Delta_\perp}
\newcommand{\FlowPot}{\Phi_\Omega}
\newcommand{\FlowPotReg}{\Phi_\Omega^{\mathrm{reg}}}
\newcommand{\CoulPot}{U_\Omega}
\newcommand{\BalanceField}{\Pi_\Omega}
\newcommand{\BalMult}{\Lambda_\Pi}
\newcommand{\SourceDensity}{\varrho_\Omega}
\newcommand{\SourceDensityRed}{\varrho_\Omega^{\mathrm{red}}}
\newcommand{\SourceMult}{\Lambda_\rho}
\newcommand{\SourceFlux}{\mathcal E}
\newcommand{\HamChargeOmega}{\mathfrak m_\Omega}
\newcommand{\SigmaIER}{\Sigma^{\mathrm{IER}}}
\newcommand{\SigmaDressSch}{\Sigma^{\mathrm{Sch}}}
\newcommand{\SigmaRegPol}{\Sigma^{\mathrm{pol,reg}}}
\newcommand{\LiftIER}{\widehat\Sigma^{\mathrm{IER}}}
\newcommand{\AuxPol}{\mathbb K}
\newcommand{\CandAction}{S_{\mathrm{cand}}}
\newcommand{\IntrinsicAction}{S_{\mathrm{cand}}^{\mathrm{src+aux}}}
\newcommand{\ReOff}{\widetilde{\mathcal R}_{\mu\nu}^{\mathrm{off}}}
\newcommand{\ReLongThree}{\widetilde{\mathcal R}_{\mu\nu}^{(\chi,3)}}

\newtheorem{definition}{Definition}
\newtheorem{proposition}{Proposition}
\newtheorem{remark}{Remark}
\newtheorem{corollary}{Corollary}
\newtheorem{theorem}{Theorem}

\title{\textbf{The \TheoryName{}}\\[0.4em]
\large Nonlocal Bridge Self-Response Intrinsic Source-Sector Note}
\author{}
\date{}

\begin{document}

\maketitle

\begin{abstract}
The positive quotient reduced-system, theorem, decay, and asymptotic line remains closed and is not reopened here. The inverse-Einstein-response bridge map, the explicit Schwarzschild-type dressing candidate test, the substrate-origin note, and the constrained auxiliary-sector note are also fixed in status: the inverse-response correction already supplies matter-free activity, full off-longitudinal \(Q/W\) cancellation, and explicit cubic longitudinal completion; the dressing-candidate note records the positive observer-level \(r^{-2}\) closure; the substrate-origin note identifies the minimal polarization mechanism \(\BalanceField=\FlowPot^2\); and the constrained auxiliary-sector note derives the polarization tensor variationally. The remaining defect is narrower: \(\FlowPot\) is still treated there as already-given input data rather than as a field produced by an intrinsic substrate source sector.

The present manuscript internalizes that still-given potential. It enlarges the same candidate substrate action by an intrinsic Hamiltonian-charge source sector consisting of an explicit substrate charge density \(\SourceDensity\), a spatial source flux \(\SourceFlux_A\), and a Lagrange multiplier \(\SourceMult\), together with the already-recorded polarization variables \((\BalanceField,\BalMult,\AuxPol_{AB})\). The new source Lagrangian yields, by Euler--Lagrange variation,
\[
\SourceDensity=\SourceDensityRed,
\qquad
\SourceFlux_A=\SpatialD_A\FlowPot,
\qquad
-\LapPerp \FlowPot=4\pi G_N\,\SourceDensity,
\]
where \(\SourceDensityRed\) is the same already-fixed reduced Hamiltonian-charge density functional whose slice integral gives \(\HamChargeOmega\). Hence \(\FlowPot\) is no longer prescribed externally: it is solved from an intrinsic substrate charge variable.

On the source-free exterior region, the unique decaying harmonic solution with total charge \(\HamChargeOmega\) has the Coulomb profile
\[
\FlowPot
=
\frac{G_N\,\HamChargeOmega}{r}
+\FlowPotReg,
\qquad
\FlowPotReg=\mathcal O(r^{-2}),
\qquad
\partial \FlowPotReg=\mathcal O(r^{-3}).
\]
The auxiliary polarization block then yields the same eliminated tensor as before,
\[
\AuxPol_{AB}^{\mathrm{elim}}
=
-2\FlowPot^2U_AU_B
+\frac32\FlowPot^2\ProjSub_{AB},
\]
so its observer pullback reproduces exactly the same Schwarzschild-type \(r^{-2}\) dressing together with an \(\mathcal O(r^{-3})\) regular completion. The enlarged source-plus-polarization sector therefore gives the same observer output that already passed the positive dressing test.

Because the source sector is elliptic and nonpropagating on each local experiential slice, and because it enters the bridge metric only through the quartic combination \(\BalanceField=\FlowPot^2\), it does not reopen the lower-order quotient PDE line. Matter-free activity survives because the same reduced Hamiltonian charge can be nonzero on the rotational matter-free regime; off-longitudinal \(Q/W\) cancellation and explicit cubic longitudinal completion survive because \(\SigmaIER\) is unchanged; and quartic Coulomb self-response absorption survives because the eliminated observer tensor is the same explicit dressed candidate already tested positively.

This note is still not a completed microphysical derivation of Einsteinian gravity from substrate first principles. Its narrower positive result is the one the active sequence now required: the Hamiltonian-charge potential \(\FlowPot\) is internalized into an enlarged action by explicit substrate charge variables, its Coulomb tail is derived from a source equation, and the same \(r^{-2}\) dressed bridge output together with the same gain package is preserved.
\end{abstract}

\section{Introduction}

The active sequence is now precise about what should happen next. The positive quotient reduced-system, theorem, decay, and asymptotic line is already recorded and remains closed at current scope. The observer-level redesign chain is also fixed in status: the inverse-Einstein-response bridge-map test already canceled the full off-longitudinal \(Q/W\) remainder and added an explicit cubic longitudinal completion, the quartic continuation note isolated the Coulomb self-response obstruction, the nonlocal self-response redesign note fixed the \(r^{-2}\) target, and the explicit dressing-candidate test recorded the positive Schwarzschild-type closure. The substrate-origin and constrained auxiliary-sector notes then took the next two reconstruction steps by identifying and variationally realizing the polarization mechanism
\[
\BalanceField=\FlowPot^2.
\]

That still leaves one real burden. In the constrained auxiliary-sector note, the polarization tensor is solved variationally, but the potential \(\FlowPot\) that sources it is still imported into the action as already-given nonlocal data. The present manuscript addresses exactly that defect and no broader one.

\paragraph{Framework and claim boundary.}
\TheoryProgram{} denotes the broader \Aether{} / \Aflow{} framework built on the claims that the \Aether{} is the underlying four-dimensional substrate of reality, the \Aflow{} is its intrinsic ordered motion, observed three-dimensional space is the local experiential slice of that deeper substrate, S-time is the experienced order of change arising from matter, light, and the \Aflow{}, and observed expansion is the three-dimensional appearance of deeper four-dimensional ordered motion. Within that framework, \TheoryName{} denotes the exact-closure theory in which the effective gravitational dynamics are adopted to be exactly Einsteinian with universal matter coupling. In the active sequence, \emph{adoption} means use of that established relativistic dynamics without claiming substrate derivation, while \emph{derivation} is reserved for a first-principles recovery from explicit substrate variables.

The present note stays strictly inside that boundary. It does not reopen the already-positive quotient PDE line. It does not replace the already-positive observer-level dressing candidate. It does not claim that the full source density functional \(\SourceDensityRed\) has already been reduced to a unique local microphysical formula. Its narrower positive result is that the potential \(\FlowPot\) can now be solved inside an enlarged substrate action from an explicit Hamiltonian-charge source variable, and that eliminating the full source-plus-polarization sector reproduces the same \(r^{-2}\) dressed observer output already fixed by the positive dressing test.

\section{Fixed Target and Present Scope}

\begin{definition}[Surviving dressed gain package]
\label{def:gains}
The \emph{surviving dressed gain package} consists of the four already-recorded features that the enlarged source-plus-polarization sector must preserve:
\begin{enumerate}
    \item \textbf{matter-free activity:} the bridge correction remains active on matter-free reduced data;
    \item \textbf{off-longitudinal \(Q/W\) cancellation:} the full currently recorded off-longitudinal remainder \(\ReOff\) is canceled;
    \item \textbf{explicit cubic longitudinal completion:} the first surviving cubic longitudinal tensor \(\ReLongThree\) is canceled by the same lower-order inverse-response correction;
    \item \textbf{quartic Coulomb self-response absorption:} on the decisive matter-free rotational exterior regime, the isolated quartic Coulomb self-response is pushed to \(\mathcal O(r^{-5})\).
\end{enumerate}
\end{definition}

\begin{definition}[Observer-side target already fixed]
\label{def:target}
The observer-side dressed bridge map to be reproduced remains
\begin{equation}
g_{\mu\nu}^{\mathrm{dressed}}
=
g_{\mu\nu}^{\mathrm{br}}
+\SigmaIER_{\mu\nu}
+\SigmaDressSch{}_{\mu\nu}
+\SigmaRegPol{}_{\mu\nu},
\label{eq:targetmetric}
\end{equation}
where the charge-carrying Coulomb piece is
\begin{equation}
\CoulPot(r)\equiv \frac{G_N\,\HamChargeOmega}{r},
\label{eq:coulpot}
\end{equation}
and
\begin{equation}
\SigmaDressSch{}_{00}=-2\CoulPot^2,
\qquad
\SigmaDressSch{}_{0i}=0,
\qquad
\SigmaDressSch{}_{ij}=\frac32\CoulPot^2\,\delta_{ij},
\label{eq:schdress}
\end{equation}
while the regular completion satisfies
\begin{equation}
\SigmaRegPol=\mathcal O(r^{-3}),
\qquad
\partial\SigmaRegPol=\mathcal O(r^{-4})
\label{eq:reg}
\end{equation}
on the exterior charge sector.
\end{definition}

\begin{remark}
The present note does not replace the already-fixed lower-order inverse-response tensor \(\SigmaIER\). Its task is to derive the Coulomb potential \(\FlowPot\) from an intrinsic source variable and then show that the resulting eliminated polarization block still reproduces Definition \ref{def:target}.
\end{remark}

\section{Intrinsic Hamiltonian Source Variables}

The substrate-kinematics manuscript already fixed the minimal substrate data
\[
(\Sub,G_{AB},U^A,S,Q),
\]
with \(G_{AB}\) positive definite, \(U^A\) the unit ordered-flow field, and
\begin{equation}
\ProjSub_{AB}=G_{AB}-U_AU_B
\label{eq:projector}
\end{equation}
the local experiential projector orthogonal to \(U^A\). The present note appends only the extra variables needed to internalize \(\FlowPot\).

\begin{definition}[Reduced Hamiltonian source density]
\label{def:redsource}
Let
\begin{equation}
\SourceDensityRed=\SourceDensityRed[Q,W,\chi]
\label{eq:redsource}
\end{equation}
denote the already-fixed reduced Hamiltonian-charge density functional on the active branch. Its exact coefficient-level formula is not reopened here. The only properties needed at the present scope are:
\begin{enumerate}
    \item \(\SourceDensityRed=\mathcal O_{\ge 2}\) in the reduced-data counting;
    \item its slice integral defines the same Hamiltonian charge already used in the inverse-response and dressing notes,
    \begin{equation}
    \HamChargeOmega
    \equiv
    \int_{\Sigma_t}\SourceDensityRed\,d\mu_P;
    \label{eq:charge}
    \end{equation}
    \item on the decisive rotational matter-free regime, \(\HamChargeOmega\) may be nonzero.
\end{enumerate}
\end{definition}

\begin{remark}
Definition \ref{def:redsource} is the precise point at which the present note stops short of a deeper microphysical claim. The source density that carries the Hamiltonian charge is now made into an explicit substrate variable, but the fully local derivation of the functional \(\SourceDensityRed[Q,W,\chi]\) itself is not yet claimed here.
\end{remark}

\begin{definition}[Intrinsic source variables]
\label{def:sourcevars}
The \emph{intrinsic Hamiltonian source sector} consists of:
\begin{enumerate}
    \item a scalar substrate charge density \(\SourceDensity\);
    \item a spatial source flux \(\SourceFlux_A\) satisfying
    \begin{equation}
    U^A\SourceFlux_A=0;
    \label{eq:spatialflux}
    \end{equation}
    \item the potential \(\FlowPot\), now treated as an independent substrate field rather than as fixed input data;
    \item a scalar Lagrange multiplier \(\SourceMult\).
\end{enumerate}
\end{definition}

\begin{definition}[Projected source derivative]
\label{def:projected}
Let
\begin{equation}
\SpatialD_A\equiv \ProjSub_A{}^B\nabla_B,
\qquad
\LapPerp\equiv \SpatialD^A\SpatialD_A,
\label{eq:projderiv}
\end{equation}
where \(\nabla_A\) is the Levi-Civita connection of \(G_{AB}\). Thus \(\SpatialD\) and \(\LapPerp\) act along the local experiential slice orthogonal to \(U^A\).
\end{definition}

\begin{definition}[Intrinsic source-plus-polarization action]
\label{def:action}
The \emph{intrinsic source-plus-polarization sector} is the enlarged action
\begin{equation}
\IntrinsicAction
\equiv
\CandAction
+
\int_{\Sub} d^4X\,\sqrt{G}\,
\bigl(
\mathcal L_{\mathrm{src}}
+\mathcal L_{\mathrm{pol,aux}}
\bigr),
\label{eq:fullaction}
\end{equation}
with source Lagrangian
\begin{equation}
\mathcal L_{\mathrm{src}}
=
\SourceFlux^A\SpatialD_A\FlowPot
-\frac12 \SourceFlux_A\SourceFlux^A
-4\pi G_N\,\SourceDensity\,\FlowPot
+\SourceMult\bigl(\SourceDensity-\SourceDensityRed\bigr),
\label{eq:lsrc}
\end{equation}
and polarization Lagrangian
\begin{align}
\mathcal L_{\mathrm{pol,aux}}
=\;&
\BalMult\bigl(\BalanceField-\FlowPot^2\bigr)
-\frac{\mu_t^2}{2}\bigl(\alpha+2\BalanceField\bigr)^2
-\frac{\mu_s^2}{2}\bigl(\beta-\tfrac32\BalanceField\bigr)^2
\nonumber\\
&-\frac{\mu_v^2}{2}\,\mathcal V_A\mathcal V^A
-\frac{\mu_{\mathrm{tf}}^2}{2}\,\mathcal T_{AB}\mathcal T^{AB},
\label{eq:laux}
\end{align}
where the independent symmetric auxiliary tensor \(\AuxPol_{AB}\) is decomposed as
\begin{equation}
\AuxPol_{AB}
=
\alpha\,U_AU_B
+2U_{(A}\mathcal V_{B)}
+\mathcal T_{AB}
+\beta\,\ProjSub_{AB},
\label{eq:Kdecomp}
\end{equation}
with
\begin{equation}
U^A\mathcal V_A=0,
\qquad
U^A\mathcal T_{AB}=0,
\qquad
\ProjSub^{AB}\mathcal T_{AB}=0,
\label{eq:TVcond}
\end{equation}
and
\begin{equation}
\mu_t^2>0,
\qquad
\mu_s^2>0,
\qquad
\mu_v^2>0,
\qquad
\mu_{\mathrm{tf}}^2>0.
\label{eq:masses}
\end{equation}
\end{definition}

\begin{remark}
Compared with the constrained auxiliary-sector note, the only conceptual change in Definition \ref{def:action} is that \(\FlowPot\) has been promoted from fixed input to a field solved from the intrinsic source variables \((\SourceDensity,\SourceFlux_A,\SourceMult)\). The already-positive lower-order candidate substrate action \(\CandAction\) and the already-positive lower-order inverse-response correction \(\SigmaIER\) are both kept fixed.
\end{remark}

\section{Variational Derivation of the Coulomb Potential}

\begin{proposition}[Euler--Lagrange system of the enlarged source-plus-polarization sector]
\label{prop:EL}
The Euler--Lagrange equations of \eqref{eq:fullaction} are
\begin{equation}
\SourceDensity=\SourceDensityRed,
\qquad
\SourceMult=4\pi G_N\,\FlowPot,
\qquad
\SourceFlux_A=\SpatialD_A\FlowPot,
\label{eq:sourceeqs}
\end{equation}
\begin{equation}
-\LapPerp \FlowPot=4\pi G_N\,\SourceDensity,
\label{eq:poisson}
\end{equation}
and
\begin{equation}
\BalanceField=\FlowPot^2,
\qquad
\alpha=-2\BalanceField,
\qquad
\beta=\frac32\BalanceField,
\qquad
\mathcal V_A=0,
\qquad
\mathcal T_{AB}=0,
\qquad
\BalMult=0.
\label{eq:poleqs}
\end{equation}
Consequently,
\begin{equation}
\AuxPol_{AB}^{\mathrm{elim}}
=
-2\FlowPot^2U_AU_B
+\frac32\FlowPot^2\ProjSub_{AB}.
\label{eq:Kelim}
\end{equation}
\end{proposition}

\begin{proof}
Variation with respect to \(\SourceMult\) gives \(\SourceDensity=\SourceDensityRed\). Variation with respect to \(\SourceDensity\) gives \(\SourceMult=4\pi G_N\FlowPot\). Variation with respect to the spatial flux \(\SourceFlux_A\) gives \(\SourceFlux_A=\SpatialD_A\FlowPot\). Varying \(\FlowPot\) and integrating by parts on the local experiential slice yields
\[
-\SpatialD_A\SourceFlux^A-4\pi G_N\,\SourceDensity-2\BalMult\FlowPot=0.
\]
The variation of \(\BalanceField\) in \eqref{eq:laux} gives
\[
\BalMult
-2\mu_t^2(\alpha+2\BalanceField)
+\frac32\mu_s^2\bigl(\beta-\tfrac32\BalanceField\bigr)=0.
\]
Variation with respect to \(\alpha\), \(\beta\), \(\mathcal V_A\), and \(\mathcal T_{AB}\) gives
\[
\alpha=-2\BalanceField,
\qquad
\beta=\frac32\BalanceField,
\qquad
\mathcal V_A=0,
\qquad
\mathcal T_{AB}=0.
\]
Substituting these into the \(\BalanceField\)-equation yields \(\BalMult=0\). The \(\FlowPot\)-equation therefore reduces to \eqref{eq:poisson}, and the multiplier equation gives \(\BalanceField=\FlowPot^2\). Substituting \eqref{eq:poleqs} into \eqref{eq:Kdecomp} yields \eqref{eq:Kelim}.
\end{proof}

\begin{corollary}[Derived exterior Coulomb profile]
\label{cor:coulomb}
Assume that \(\SourceDensityRed\) is supported in a bounded charge core \(r\le R_\ast\). Then on the source-free exterior region \(r>R_\ast\),
\begin{equation}
\FlowPot
=
\CoulPot+\FlowPotReg,
\qquad
\CoulPot(r)=\frac{G_N\,\HamChargeOmega}{r},
\qquad
\FlowPotReg=\mathcal O(r^{-2}),
\qquad
\partial\FlowPotReg=\mathcal O(r^{-3}).
\label{eq:potdecomp}
\end{equation}
\end{corollary}

\begin{proof}
By Proposition \ref{prop:EL}, the exterior region satisfies
\[
\LapPerp\FlowPot=0
\qquad
\text{for } r>R_\ast.
\]
Standard harmonic expansion on the asymptotically Euclidean local experiential slice gives
\[
\FlowPot=\frac{A}{r}+\mathcal O(r^{-2})
\]
for the unique decaying branch. Integrating \eqref{eq:poisson} over a ball \(B_r\) containing the charge core and applying the divergence theorem yields
\begin{equation}
-\int_{S_r} n^A\SpatialD_A\FlowPot\,dA
=
4\pi G_N\int_{B_r}\SourceDensity\,d\mu_P
=
4\pi G_N\,\HamChargeOmega.
\label{eq:gauss}
\end{equation}
For the harmonic monopole \(A/r\), the left-hand side of \eqref{eq:gauss} is \(4\pi A\). Hence \(A=G_N\HamChargeOmega\), which gives \eqref{eq:potdecomp}.
\end{proof}

\begin{remark}
Corollary \ref{cor:coulomb} is the precise sense in which the enlarged action now derives the Coulomb profile rather than assuming it. The charge-carrying \(r^{-1}\) tail is fixed by the intrinsic source variable \(\SourceDensity\) through Gauss law, not inserted beforehand as external potential data.
\end{remark}

\section{Elimination and Recovery of the Same \(r^{-2}\) Dressing}

\begin{definition}[Dressed bridge reconstruction after full elimination]
\label{def:bridge}
Let \(\LiftIER_{AB}\) be a substrate tensor whose observer pullback is the already-fixed lower-order inverse-response correction,
\begin{equation}
\ObserverMap^\ast\LiftIER_{AB}=\SigmaIER_{\mu\nu}.
\label{eq:liftier}
\end{equation}
After eliminating the intrinsic source-plus-polarization sector, define the reconstructed bridge metric by
\begin{equation}
\BridgeMetric_{AB}^{\mathrm{src+aux}}
=
G_{AB}
-2U_AU_B
+\LiftIER_{AB}
+\AuxPol_{AB}^{\mathrm{elim}}.
\label{eq:bridge}
\end{equation}
\end{definition}

\begin{proposition}[Coulomb-plus-regular split of the eliminated polarization tensor]
\label{prop:split}
On the exterior charge sector,
\begin{equation}
\AuxPol_{AB}^{\mathrm{elim}}
=
\AuxPol_{AB}^{\mathrm{Sch}}
+\AuxPol_{AB}^{\mathrm{reg}},
\label{eq:Ksplit}
\end{equation}
where
\begin{equation}
\AuxPol_{AB}^{\mathrm{Sch}}
=
-2\CoulPot^2U_AU_B
+\frac32\CoulPot^2\ProjSub_{AB},
\label{eq:Ksch}
\end{equation}
and
\begin{align}
\AuxPol_{AB}^{\mathrm{reg}}
=\;&
-2\bigl(2\CoulPot\FlowPotReg+(\FlowPotReg)^2\bigr)U_AU_B
\nonumber\\
&+\frac32\bigl(2\CoulPot\FlowPotReg+(\FlowPotReg)^2\bigr)\ProjSub_{AB}.
\label{eq:Kreg}
\end{align}
Moreover,
\begin{equation}
\AuxPol_{AB}^{\mathrm{reg}}=\mathcal O(r^{-3}),
\qquad
\nabla\AuxPol_{AB}^{\mathrm{reg}}=\mathcal O(r^{-4}).
\label{eq:Kregdecay}
\end{equation}
\end{proposition}

\begin{proof}
Substitute \eqref{eq:potdecomp} into \eqref{eq:Kelim}:
\[
\FlowPot^2=\CoulPot^2+2\CoulPot\FlowPotReg+(\FlowPotReg)^2.
\]
Separating the pure Coulomb square from the faster remainder gives \eqref{eq:Ksplit}--\eqref{eq:Kreg}. Since \(\CoulPot=\mathcal O(r^{-1})\) and \(\FlowPotReg=\mathcal O(r^{-2})\), one has \(2\CoulPot\FlowPotReg+(\FlowPotReg)^2=\mathcal O(r^{-3})\), and the derivative estimate follows immediately.
\end{proof}

\begin{proposition}[Observer pullback of the full eliminated source-plus-polarization sector]
\label{prop:pullback}
In a substrate-adapted local observer chart in which \(U^A\) is the unit temporal direction and \(\ObserverMap^\ast\ProjSub_{AB}\) reduces to the spatial Euclidean metric at leading order, one has
\begin{equation}
\bigl(\ObserverMap^\ast\AuxPol^{\mathrm{Sch}}\bigr)_{00}
=
-2\CoulPot^2,
\qquad
\bigl(\ObserverMap^\ast\AuxPol^{\mathrm{Sch}}\bigr)_{0i}
=
0,
\qquad
\bigl(\ObserverMap^\ast\AuxPol^{\mathrm{Sch}}\bigr)_{ij}
=
\frac32\CoulPot^2\,\delta_{ij},
\label{eq:pullsch}
\end{equation}
and
\begin{equation}
\ObserverMap^\ast\AuxPol_{AB}^{\mathrm{reg}}
=
\SigmaRegPol{}_{\mu\nu},
\qquad
\SigmaRegPol=\mathcal O(r^{-3}),
\qquad
\partial\SigmaRegPol=\mathcal O(r^{-4}).
\label{eq:pullreg}
\end{equation}
Therefore
\begin{equation}
\ObserverMap^\ast \BridgeMetric_{AB}^{\mathrm{src+aux}}
=
g_{\mu\nu}^{\mathrm{br}}
+\SigmaIER_{\mu\nu}
+\SigmaDressSch{}_{\mu\nu}
+\SigmaRegPol{}_{\mu\nu},
\label{eq:targetrecovered}
\end{equation}
with \(\SigmaDressSch\) exactly as in \eqref{eq:schdress}.
\end{proposition}

\begin{proof}
The pullback identities for \(U_AU_B\) and \(\ProjSub_{AB}\) are the same ones used in the substrate-origin and constrained auxiliary-sector notes. Applying them to \eqref{eq:Ksch} yields \eqref{eq:pullsch}, which is exactly the Schwarzschild-type \(r^{-2}\) dressing of Definition \ref{def:target}. Applying the same pullback to \eqref{eq:Kreg} and using \eqref{eq:Kregdecay} gives \eqref{eq:pullreg}. Combining those identities with \eqref{eq:bridge} and \eqref{eq:liftier} yields \eqref{eq:targetrecovered}.
\end{proof}

\begin{corollary}[Same observer output as the positive dressing test]
\label{cor:sameoutput}
Eliminating the intrinsic source-plus-polarization sector reproduces exactly the same observer-side dressed bridge map already fixed by the positive dressing-candidate test.
\end{corollary}

\begin{proof}
This is precisely Proposition \ref{prop:pullback}.
\end{proof}

\section{Why the Same Gain Package Survives}

\begin{proposition}[Source sector is elliptic and does not introduce new propagating observer modes]
\label{prop:elliptic}
At the present manuscript scope, the intrinsic Hamiltonian source sector is nonpropagating on each local experiential slice:
\begin{equation}
\SourceDensity=\mathcal O_{\ge 2},
\qquad
\FlowPot=\mathcal O_{\ge 2},
\qquad
\BalanceField=\FlowPot^2=\mathcal O_{\ge 4}.
\label{eq:orders}
\end{equation}
Moreover, the source sector enters the bridge metric only through the quartic combination \(\BalanceField=\FlowPot^2\).
\end{proposition}

\begin{proof}
Definition \ref{def:redsource} gives \(\SourceDensityRed=\mathcal O_{\ge 2}\), hence Proposition \ref{prop:EL} gives \(\SourceDensity=\mathcal O_{\ge 2}\). The elliptic solve \eqref{eq:poisson} is linear in \(\SourceDensity\), so \(\FlowPot=\mathcal O_{\ge 2}\) in the same reduced-data counting used in the inverse-response and dressing notes. Therefore \(\BalanceField=\FlowPot^2=\mathcal O_{\ge 4}\). Equation \eqref{eq:Kelim} shows that the source sector reaches the bridge metric only through \(\BalanceField\), hence only at quartic order.
\end{proof}

\begin{corollary}[No reopening of the lower-order quotient line]
\label{cor:nopde}
The intrinsic source sector does not reopen the already-positive quotient reduced-system, theorem, decay, or asymptotic line.
\end{corollary}

\begin{proof}
By Proposition \ref{prop:elliptic}, the new source variables add only a slice-wise elliptic reconstruction whose bridge imprint begins at quartic order. Therefore they do not alter the already-closed lower-order quotient equations or the estimates built on them.
\end{proof}

\begin{theorem}[Preservation of matter-free activity, off-longitudinal cancellation, and cubic completion]
\label{thm:firstthree}
The intrinsic source-plus-polarization sector preserves items (1)--(3) of Definition \ref{def:gains}: matter-free activity, off-longitudinal \(Q/W\) cancellation, and explicit cubic longitudinal completion.
\end{theorem}

\begin{proof}
The lower-order inverse-response tensor \(\SigmaIER\) is unchanged by Definition \ref{def:bridge}. Therefore the lower-order identities through which \(\SigmaIER\) cancels the full off-longitudinal remainder \(\ReOff\) and the cubic longitudinal tensor \(\ReLongThree\) remain exactly as before. Matter-free activity also survives: by Definition \ref{def:redsource}, the same reduced Hamiltonian charge \(\HamChargeOmega\) may be nonzero on the decisive rotational matter-free regime, so the elliptic solve \eqref{eq:poisson} produces the same nontrivial Coulomb tail even when ordinary matter vanishes. Since Proposition \ref{prop:elliptic} shows that the source-plus-polarization sector enters the observer map only at quartic order, none of these already-positive lower-order gains is disturbed.
\end{proof}

\begin{theorem}[Preservation of quartic Coulomb self-response absorption]
\label{thm:quartic}
On the decisive matter-free rotational exterior regime, the intrinsic source-plus-polarization sector preserves item (4) of Definition \ref{def:gains}: the isolated quartic Coulomb self-response is structurally absorbed and the resulting quartic observer tensor is pushed to \(\mathcal O(r^{-5})\).
\end{theorem}

\begin{proof}
By Proposition \ref{prop:pullback}, the observer tensor produced by eliminating the full source-plus-polarization sector is exactly the same Schwarzschild-type \(r^{-2}\) dressing plus an \(\mathcal O(r^{-3})\) regular completion already tested positively in the dressing-candidate note. The Hamiltonian projection cancellation and full tensor-level \(\mathcal O(r^{-5})\) verdict of that note therefore apply verbatim.
\end{proof}

\begin{corollary}[Intrinsic source-sector verdict]
\label{cor:verdict}
At the present disciplined scope, the enlarged intrinsic source-plus-polarization sector of Definition \ref{def:action} is a valid next-step realization of the surviving dressed bridge map:
\begin{enumerate}
    \item it introduces an explicit substrate Hamiltonian charge variable \(\SourceDensity\) for the source of \(\FlowPot\);
    \item it derives \(\FlowPot\) from the source equation \(-\LapPerp\FlowPot=4\pi G_N\SourceDensity\) instead of taking \(\FlowPot\) as fixed input;
    \item it reproduces the same Schwarzschild-type \(r^{-2}\) dressing and \(\mathcal O(r^{-3})\) regular completion after eliminating the full source-plus-polarization sector;
    \item it preserves the same matter-free/off-longitudinal/cubic/quartic gain package.
\end{enumerate}
\end{corollary}

\begin{proof}
Item (1) is Definition \ref{def:sourcevars}. Item (2) is Proposition \ref{prop:EL} together with Corollary \ref{cor:coulomb}. Item (3) is Proposition \ref{prop:pullback}. Item (4) is Theorems \ref{thm:firstthree} and \ref{thm:quartic}.
\end{proof}

\section{Conclusion}

The active burden at this stage was no longer another observer-level repair, and it was no longer merely the constrained auxiliary-sector reconstruction either. It was to stop the Hamiltonian-charge potential \(\FlowPot\) from remaining an input. The present manuscript supplies that next step.

\begin{enumerate}
    \item The enlarged action now contains an explicit substrate Hamiltonian charge density \(\SourceDensity\), a spatial source flux \(\SourceFlux_A\), and a multiplier \(\SourceMult\), so the source of \(\FlowPot\) is now represented by explicit substrate variables.
    \item The Euler--Lagrange equations derive the slice-wise Poisson equation \(-\LapPerp\FlowPot=4\pi G_N\SourceDensity\), so the Coulomb profile is now produced from the source sector rather than prescribed beforehand.
    \item The same auxiliary polarization block still yields \(\BalanceField=\FlowPot^2\) and \(\AuxPol_{AB}^{\mathrm{elim}}=-2\FlowPot^2U_AU_B+\tfrac32\FlowPot^2\ProjSub_{AB}\).
    \item On the exterior charge sector, that elimination splits into the same Schwarzschild-type \(r^{-2}\) dressing plus an \(\mathcal O(r^{-3})\) regular completion, so the observer output is exactly the same dressed bridge map already fixed by the positive candidate test.
    \item Because the new source sector is elliptic and enters the bridge only through the quartic combination \(\FlowPot^2\), the already-positive matter-free activity, off-longitudinal \(Q/W\) cancellation, explicit cubic longitudinal completion, and quartic Coulomb self-response absorption are all preserved.
\end{enumerate}

This is the correct scope for the present note. It does not yet claim that the reduced Hamiltonian source density \(\SourceDensityRed[Q,W,\chi]\) has been derived from a unique deeper local microphysics, and it does not yet promote the full enlarged bridge variable into a deeper local substrate curvature action. The narrower positive result is the one the active sequence now required: the still-given potential \(\FlowPot\) has been internalized into an intrinsic substrate source sector without losing the same \(r^{-2}\) dressed bridge output or the same gain package.

\end{document}
